\section{Related work}
\label{sec:related}
\textbf{Exposure mappings and design-based estimation.}
Design-based work defines the effect of an \emph{exposure} (a function of the unit's own and its neighbours' treatments) and estimates
contrasts of exposures by inverse-probability weighting, with variance bounds that hold for arbitrary interference once the exposure mapping is
correct \citep{aronowsamii2017}. The same ideas, with graph-cluster randomisation to make ``all neighbours treated'' likely, give Horvitz--Thompson
estimators whose variance grows with degree under independent assignment but can be controlled by clustering on restricted-growth graphs
\citep{ugander2013}; \citet{eckles2017} show that cluster randomisation and neighbour-aware analysis reduce the bias of GATE estimation in simulations,
most when the network is clustered and interference is strong, and \citet{karrer2021} describe an industrial deployment. \citet{chen2023} derive
bias and variance of an HT estimator under a one-hop model as a function of topology and the covariance of the assignment, and optimise the latter.
Our Hájek baselines follow this line; their failure when the exposure model is wrong is documented in these works and in \citet{chin2019}.

\textbf{Regression-based GATE estimation.}
\citet{chin2019} introduced the estimator we study: fit, by (separate-arm) least squares, a model of the outcome on features of the assignment of
\emph{other} units (e.g.\ the treated fraction and count of neighbours), predict outcomes at the all-treated and all-control feature values, and difference.
Under Bernoulli assignment, strict exogeneity and a correctly specified linear model it is unbiased with an exact variance expression whose
extrapolation term makes it several times more variable than the difference in means, and it comes with a distributional result that need not
be $\sqrt n$ under interference; flexible (additive-model) fits were reported to add little over OLS in one nonlinear simulation, and bootstrap
intervals rely on exogeneity.
\citet{hanugander2023} select graph- and treatment-based covariates from a larger generated set and give block-bootstrap intervals; their
guarantee is consistency (sparse linear truth, restricted-eigenvalue condition) rather than a finite-sample error expression, and they report that
simple fraction-only adjustment can be badly biased under misspecification and that naive bootstrap coverage can be well below nominal.
\citet{cortez2022} take a different route: under staggered rollout, GATE is estimable without network knowledge when interference is a low-degree
polynomial of the neighbours' treatments, which can be read as polynomial extrapolation of the same kind we study; its variance grows steeply in the degree of the polynomial.
We did not implement staggered rollout (Sec.~\ref{sec:limits}).

\textbf{What is new here.}
The cited works characterise the estimator under correct specification, prove consistency, or simulate particular networks and outcome models, and \citet{chin2019} and \citet{hanugander2023} already report that fraction-only adjustment can be biased under misspecification. We do not claim a new principle. We add (i)~a closed form of that bias for the neighbour-fraction OLS under one-hop spillover and Bernoulli assignment, evaluated on the actual degree mixture and tested against simulations with exactly known truth; (ii)~a controlled, quantitative map of error, coverage and baselines across graphs, mechanisms, designs and specifications, with bias and variance reported separately and a correctly specified control. Our positioning rests on the works cited here; we claim no priority beyond them.
